% GPT-6 Astra Pro; effort not exposed; /root/replacement; 22 November 2026.
\section[HD-F1-R29: rank-one factors force abelian supports]{HD-F1-R29: rank-one divisor factors force abelian surface supports and are excluded in the three-row branch}
\label{sec:followup-rm-rank-one-surface-support}

Retain the actual objects and the complete product
characters of R28. In the three-row common-image
branch, both semidefinite divisor factors have
rank two. Thus both are definite, with the exact
signs already specified in R28. In the
common-quotient branch, rank-one factors require
$r\ne0$ and, in the explicit reparametrization
(F1R28.6),
\[
       \tau=\overline\tau=\frac{a_0q}{t},\qquad
       t=2-a_0>0.
 \tag{F1R29.1}
\]
Moreover every component of its actual target
surface defect is a translate of the same
abelian surface, with one tangent direction in
each RM block. These are restrictions on the
original sheaves and their actual quotient
cycles. The remaining definite-factor cases
and the two-row common-image branch are retained.

\paragraph{1. An effective surface with a single rank-one mixed factor.}
Let $T_1\oplus T_2=T_0X$ be the genuine RM
decomposition, and let $\mu_i$ be nonzero
semipositive invariant $(1,1)$ forms of rank one
on $T_i$, extended by zero on the other block.
They need not be rational divisor classes.
For every effective surface cycle $D$ whose
actual cohomology class satisfies
\[
                 [D]=uA_1+vA_2+w\mu_1\mu_2,
                          \qquad u,v,w\in\mathbf R,
 \tag{F1R29.2}
\]
one has $u=v=0$. For $D\ne0$, every irreducible
component is a translate of one abelian surface
$B\subset X$, and
\[
          T_0B=\ker(\mu_1|_{T_1})
                       \oplus\ker(\mu_2|_{T_2}).
 \tag{F1R29.3}
\]
In particular $w>0$. The zero cycle has
$u=v=w=0$ by the exterior block decomposition.
We prove the nonzero assertion with the precise
duality maps, rather than identifying a form's
active direction with a tangent direction.

Choose holomorphic bases $(\alpha_i,\eta_i)$ of
$T_i^*$ so that $\mu_i=\sqrt{-1}\alpha_i\wedge
\overline\alpha_i$; absorb its positive
coefficient into $\alpha_i$. Let $(e_1,e_2)$
and $(f_1,f_2)$ be the dual tangent bases for
the first and second blocks. Thus
$\ker\mu_1=\mathbf C e_2$ and
$\ker\mu_2=\mathbf C f_2$.
Every mixed holomorphic two-form is uniquely
\[
 \omega=b_{11}\alpha_1\alpha_2
       +b_{12}\alpha_1\eta_2
       +b_{21}\eta_1\alpha_2
       +b_{22}\eta_1\eta_2.
 \tag{F1R29.4}
\]
Here and below adjacent differential forms are
wedged. Excess holomorphic degree in either
block gives $A_i\omega\overline\omega=0$.
In $\mu_1\mu_2\omega\overline\omega$, every
term vanishes except the term with coefficient
$|b_{22}|^2$. Its coefficient relative to the
positive four-dimensional product volume is
positive. Consequently the cohomological
pairing in (F1R29.2) vanishes for each of
\[
       \alpha_1\alpha_2,\qquad
       \alpha_1\eta_2,\qquad
       \eta_1\alpha_2.
 \tag{F1R29.5}
\]

Write $D=\sum n_j Z_j$, with $n_j>0$, and take
a smooth projective resolution
$g_j:Y_j\to Z_j\hookrightarrow X$.
The actual cycle-class pairing is
\[
 \int_X[D]\,\omega\overline\omega
   =\sum_j n_j\int_{Y_j}g_j^*\omega\,
                                  \overline{g_j^*\omega}.
 \tag{F1R29.6}
\]
Each summand is the integral of a nonnegative
top-degree form on the surface. Zero integral
forces the holomorphic two-form to vanish
identically. Thus all three pullbacks in
(F1R29.5) vanish on each $Y_j$.

For a smooth tangent plane of $Z_j$, evaluation
of the forms in (F1R29.5) is the pairing of
their covectors with its Pl\"ucker bivector.
They kill respectively its $e_1f_1$, $e_1f_2$
and $e_2f_1$ coefficients. The remaining
bivector is therefore exactly
\[
        \zeta=x\,e_1e_2+y\,e_2f_2+z\,f_1f_2,
       \qquad
        \zeta^2=2xz\,e_1e_2f_1f_2.
 \tag{F1R29.7}
\]
Decomposability gives $xz=0$. The two resulting
projective lines are closed in the Grassmannian.
The Gauss image of the irreducible smooth locus
has irreducible closure, so it lies in one of
these lines. In the line $z=0$ every plane is
contained in $T_1\oplus\mathbf C f_2$ and is
annihilated by the fixed form $\alpha_2$.
In the line $x=0$ it is contained in
$\mathbf C e_2\oplus T_2$ and is annihilated
by the fixed form $\alpha_1$.
That one-form consequently pulls back to zero
on the entire resolution.

\paragraph{2. The actual Albanese image and the irrational quotient.}
Translate $g_j$ so that a chosen base point maps
to zero. Its induced Albanese homomorphism
\[
              \phi_j:\operatorname{Alb}(Y_j)
                                     \longrightarrow X
 \tag{F1R29.8}
\]
can be described directly on periods.
The dual of
$g_j^*:H^0(X,\Omega_X^1)\to H^0(Y_j,\Omega_{Y_j}^1)$
sends the period of a loop $\gamma$ to the
period of $(g_j)_*\gamma$. Hence it sends
the source integral period lattice into the
target lattice, and descends to (F1R29.8).
Path integration from the base point gives
$\phi_j\circ\operatorname{alb}_{Y_j}=g_j$.
Its differential is exactly the dual of $g_j^*$.
This proves the tangent assertion used here
from the actual homomorphism.

Let $B_j=\operatorname{im}\phi_j$, a connected
abelian subvariety. Its dimension is at least
two because it contains the surface $Z_j$
after the chosen translation, and at most
three because one of the nonzero forms
$\alpha_i$ annihilates its tangent space.
Dimension three is impossible. In that case
\[
                         T_0B_j=\ker\alpha_i.
 \tag{F1R29.9}
\]
Choose the actual nonrational RM endomorphism
$j=m(1+\sqrt{\Delta_0})\in\operatorname{End}(X)$
from (F1R19.11). Its differential has the
two distinct irrational real eigenvalues
$\lambda_1=m(1+\sqrt{\Delta_0})$ and
$\lambda_2=m(1-\sqrt{\Delta_0})$.
The hyperplane in (F1R29.9) is invariant under
this differential. Therefore $j(B_j)=B_j$:
both connected images have exactly that Lie
algebra, and a connected subtorus is the image
of its Lie algebra in the universal cover.
The quotient homomorphism
\[
       \pi:X\longrightarrow X/B_j,
       \qquad \pi j=\overline j\,\pi
 \tag{F1R29.10}
\]
then gives an endomorphism $\overline j$ of the
one-dimensional quotient torus. Explicitly,
if $X=V/\Lambda$ and $W=T_0B_j$, this quotient
is $(V/W)/\operatorname{im}(\Lambda\to V/W)$;
the displayed lattice is identified with
$\Lambda/(\Lambda\cap W)$.
Its differential is multiplication
by the irrational real number $\lambda_i$.
On its underlying real homology the same
map is scalar multiplication by $\lambda_i$,
but it preserves the integral lattice.
A nonzero rational coordinate of any nonzero
rational homology vector would therefore
express $\lambda_i$ as a ratio of two rational
coordinates. This contradiction rules out
dimension three.

Thus $\dim B_j=2$. The irreducible surface
$Z_j$ is a closed subvariety of a translate
of $B_j$ of the same dimension. It equals
that translate. In particular its tangent
plane $W_j=T_0B_j$ is constant.

\paragraph{3. An actual RM intersection removes both pure tangent coefficients.}
We next prove that this constant plane is
exactly the plane in (F1R29.3).
Consider the line $z=0$ in (F1R29.7).
If $x\ne0$ and $y=0$, its plane is $T_1$.
This cannot be an abelian-subvariety tangent
space: the inclusion gives nonzero rational
homology in the underlying real $T_1$
eigenspace, where the rational endomorphism
$j_*$ acts by the irrational scalar
$\lambda_1$. The same rational-coordinate
argument as above is a contradiction.

If $x\ne0$ and $y\ne0$, the exact tangent
planes of the actual abelian subvarieties
$B_j$ and $j(B_j)$ are, for some $a\ne0$,
\[
 \begin{split}
 W_j&=\langle e_2,e_1+a f_2\rangle,\\
 dj(W_j)&=\langle e_2,
                      \lambda_1e_1+a\lambda_2 f_2\rangle,\\
 W_j\cap dj(W_j)&=\mathbf C e_2.
 \end{split}
 \tag{F1R29.11}
\]
The last equality uses $a\ne0$ and
$\lambda_1\ne\lambda_2$.
The identity component
$C_j=(B_j\cap j(B_j))^0$ is an actual abelian
subvariety with this one-dimensional tangent
space. The tangent identity follows directly
by lifting a sufficiently small neighborhood
of zero to the universal cover $V$: the two
subgroups lift there to their linear tangent
subspaces, and their intersection lifts to
the intersection of those subspaces.
Its nonzero rational homology would now lie
in the underlying real $\mathbf C e_2\subset
T_1$, again contradicting the irrational
$j_*$ eigenvalue. Thus $x=0$.
For the other projective line the same
calculation, with the two blocks exchanged,
gives $z=0$.

Every component consequently has tangent
space $\mathbf C e_2\oplus\mathbf C f_2$.
Connected subtori with the same tangent
space coincide, by their universal-cover
images. Hence all $B_j$ are one $B$,
proving (F1R29.3). That tangent space is
RM-invariant. The actual endomorphisms
from a denominator-cleared order in $F$
preserve $B$ and restrict to it with
multiplicities $(1,1)$.

Finally the pullbacks of both $A_i$ to
each translate of $B$ are zero: its
projection to each $T_i$ has complex
dimension one, whereas $A_i$ has
holomorphic degree two in that block.
Use the actual integration pairings
\[
 0=\int_D A_1=v\int_X A_1A_2,
       \qquad
 0=\int_D A_2=u\int_X A_1A_2.
 \tag{F1R29.12}
\]
Here $\mu_1\mu_2 A_i=0$ and $A_i^2=0$,
while $\int_XA_1A_2>0$.
Thus $u=v=0$. Pairing the remaining
$w\mu_1\mu_2$ with any ample square
gives $w>0$ when $D\ne0$.
This proves paragraph~1, including its
precise cohomology-to-support conclusion.

\paragraph{4. The two actual defects in the three-row branch.}
Return to the common-image case $h=3$.
Suppose the two nonzero semidefinite
divisor factors of R28 had rank one.
Their squares are then zero.
For $r\ne0$, the exact rescaling gives
$\Gamma=\gamma+\overline\gamma$ and
$3\ell=t\Gamma$, so $\ell^2$ is purely
mixed. For $r=0$, the actual first-Chern
identity gives $3\ell=0$, and the same
assertion holds.

Use the original coherent sequences of
(F1R23.20), not only their hull classes.
Their actual effective surface cycles are
\[
 \begin{split}
 D_A&=[A^{**}/A]_2
                =\tfrac32\ell^2-\operatorname{ch}_2(A),\\
 D_T&=[I^{**}/I]_2+[Q^{**}/Q]_2
                =\tfrac52\ell^2-\operatorname{ch}_2(T).
 \end{split}
 \tag{F1R29.13}
\]
Each class is a real combination of
$A_1,A_2$ and the single product
$\beta\overline\beta$. Change the signs
of its two rank-one factors only in
the auxiliary real-form notation to
write that mixed line as
$\mathbf R\mu_1\mu_2$ with both $\mu_i$
semipositive. No sheaf is changed.
Paragraph~1 forces the pure-block part
of each actual cycle to be zero.
It also applies when one of the cycles
is zero; no nonvanishing premise for
the source defect is needed at this step.

On the other hand the complete character
difference is $\operatorname{ch}(A)-
\operatorname{ch}(T)=v_q$, whose
degree-two part is $-q(A_1+A_2)$.
Subtracting the exact identities in
(F1R29.13) gives
\[
                   D_A-D_T=-\ell^2+q(A_1+A_2).
 \tag{F1R29.14}
\]
Its pure-block part is the nonzero
$q(A_1+A_2)$, contradicting the two
vanishing pure-block parts just proved.
The rank-one case is excluded. By R28
both ranks must therefore be two,
proving the first assertion of this section.

\paragraph{5. The exact remaining rank-one common-quotient character.}
In the common-quotient branch use the
actual hull injection $T\hookrightarrow
L^{\oplus2}$ of R22. Its nonzero effective
surface cycle is
\[
             D_T=[L^{\oplus2}/T]_2
                         =\ell^2-\operatorname{ch}_2(T).
 \tag{F1R29.15}
\]
Suppose its two divisor factors have
rank one. If $r=0$, then $a_0=2$
and $2\ell=c_1(T)=0$. The actual
degree-two character formula gives
\[
                 [D_T]=2q(A_1+A_2)-c\beta\overline\beta,
 \tag{F1R29.16}
\]
contradicting paragraph~1.
Thus $r\ne0$. Use the exact tensor
identity (F1R28.6), so
$\gamma^2=\overline\gamma^2=0$,
$2\ell=t(\gamma+\overline\gamma)$ and
\[
 [D_T]=(a_0q-t\tau)A_1
       +(a_0q-t\overline\tau)A_2
                   -\frac{ta_0}{2}\gamma\overline\gamma.
 \tag{F1R29.17}
\]
The pure-block coefficients must both
vanish. This is exactly (F1R29.1).
Since $\tau\in F$ is fixed by its
nontrivial involution, the common value
in that equation is rational.
Every component of this very $D_T$
is a translate of the abelian surface
with tangent space
$\ker\gamma\oplus\ker\overline\gamma$
by paragraph~1. The coefficient is
consistent with its positive effective
orientation because R28 gives $t,a_0>0$
and opposite signs for the two factors.
This description does not assert that
the original quotient sheaf is the
structure sheaf of that support, or
that any such surface modification
has the required full trace.

\paragraph{Framework and precise source scope.}
The new maps are pullback of actual
mixed two-forms to component resolutions,
the period-lattice construction of
(F1R29.8), the original RM endomorphism
and its actual quotient and subgroup
intersection maps, and the unchanged
coherent quotient sequences in
(F1R29.13) and (F1R29.15).
The period construction is the one
in Bloch, \emph{Lectures on Algebraic
Cycles}, second edition, Lecture~1,
printed pages~9--10. Its Proposition~1.1(ii)
states the universal factorization;
the exact dual-period map needed here
is written and verified above, since
the published text leaves that proof
as an exercise. The genuine RM
rational-homology obstruction is
the one supplied and proved in
DEG\_L S-22, at its effective-surface
proof, with the new hyperplane quotient
and elliptic intersection spelled out.
No theorem about a merely degenerate
Gauss map, no classification of
subvarieties of a torus, and no new
sheaf-existence result is a premise.

